\documentclass[11pt]{article}
\input{../../shared/project_style.tex}
\rhead{Basics 3 Textbook Draft}

\begin{document}
\DocTitle{Basics 3: Vector Calculus for Plasma Physics}{I - Mathematical and computational foundations}{Gradient, divergence, curl, Laplacian, integral theorems, and physical interpretation}

\begin{overviewbox}
\textbf{Textbook draft, version 1.0.} Vector calculus turns spatial variation into physical statements: gradients drive forces, divergence measures sources and compression, and curl measures circulation and current. This chapter is designed for a reader with no prior plasma-physics training. It distinguishes exact identities from physical approximations and connects the central equations to later simulation modules.
\end{overviewbox}

\ProjectTOC

\section{Learning objectives}
After completing this note, the reader should be able to:
\begin{itemize}[leftmargin=1.6em]
\item Compute and interpret gradient, divergence, curl, and Laplacian operators.
\item Derive local conservation laws from integral balances.
\item Use the divergence theorem and Stokes theorem.
\item Apply vector identities used in MHD force balance and induction.
\item Recognize geometric factors in cylindrical coordinates.
\item Translate continuous operators into conservative numerical forms.
\end{itemize}

\section{Prerequisites and reading strategy}
\begin{itemize}[leftmargin=1.6em]
\item Basics 2 and multivariable differentiation and integration.
\end{itemize}
On a first reading, emphasize physical meaning and dimensional checks. On a second reading, reproduce the derivations. Attempt the computational laboratory only after the limiting cases are understood.

\section{Gradient}

For a scalar field $f$,
\begin{equation}
df=\nabla f\cdot d\mathbf r.
\end{equation}
The gradient points toward steepest increase. Pressure force density is $-\nabla p$, so it points toward decreasing pressure. A vector tangent to a level surface $f=\mathrm{constant}$ obeys $\mathbf t\cdot\nabla f=0$.

In Cartesian coordinates,
\begin{equation}
\nabla f=\mathbf e_x f_x+\mathbf e_y f_y+\mathbf e_z f_z.
\end{equation}
The directional derivative along unit vector $\mathbf n$ is $\mathbf n\cdot\nabla f$.


\section{Divergence and conservation}

Divergence is net outward flux per unit volume:
\begin{equation}
\nabla\cdot\mathbf F=
\lim_{V\to0}\frac{1}{V}\oint_{\partial V}\mathbf F\cdot d\mathbf S.
\end{equation}
The divergence theorem,
\begin{equation}
\int_V\nabla\cdot\mathbf F\,dV
=\oint_{\partial V}\mathbf F\cdot d\mathbf S,
\end{equation}
converts an integral balance into a local law. If $q$ is a conserved density with flux $\mathbf F$ and source $S$,
\begin{equation}
\pdv{q}{t}+\nabla\cdot\mathbf F=S.
\end{equation}
This structure is the mathematical basis of finite-volume methods.


\section{Curl and circulation}

Stokes theorem states
\begin{equation}
\int_S(\nabla\times\mathbf F)\cdot d\mathbf S
=\oint_{\partial S}\mathbf F\cdot d\boldsymbol\ell.
\end{equation}
Curl is therefore circulation per unit area. Maxwell equations use curl to connect current to magnetic field and changing magnetic flux to electric circulation:
\begin{equation}
\nabla\times\mathbf B\approx\mu_0\mathbf J,\qquad
\nabla\times\mathbf E=-\pdv{\mathbf B}{t}.
\end{equation}
Fluid vorticity is $\boldsymbol\omega=\nabla\times\mathbf u$.


\section{Laplacian and diffusion}

The scalar Laplacian is
\begin{equation}
\nabla^2f=\nabla\cdot\nabla f.
\end{equation}
In diffusion,
\begin{equation}
\pdv{f}{t}=D\nabla^2f,
\end{equation}
a local maximum generally has negative Laplacian and decreases, while a local minimum increases. The operator smooths spatial variation. Vector Laplacians require care in curvilinear coordinates because basis vectors also vary.


\section{Magnetic pressure and tension identity}

A key identity is
\begin{equation}
(\nabla\times\mathbf B)\times\mathbf B
=(\mathbf B\cdot\nabla)\mathbf B
-\nabla\left(\frac{B^2}{2}\right)
+\mathbf B(\nabla\cdot\mathbf B).
\end{equation}
With $\nabla\cdot\mathbf B=0$,
\begin{equation}
\frac{1}{\mu_0}(\nabla\times\mathbf B)\times\mathbf B
=\frac{1}{\mu_0}(\mathbf B\cdot\nabla)\mathbf B
-\nabla\left(\frac{B^2}{2\mu_0}\right).
\end{equation}
The first term is tension along curved field lines; the second is magnetic-pressure force.


\section{Induction identity and geometric operators}

Another central identity is
\begin{align}
\nabla\times(\mathbf u\times\mathbf B)
={}&\mathbf u(\nabla\cdot\mathbf B)-\mathbf B(\nabla\cdot\mathbf u)\\
&+(\mathbf B\cdot\nabla)\mathbf u-(\mathbf u\cdot\nabla)\mathbf B.
\end{align}
It separates compression, stretching, and advection.

For cylindrical $\mathbf A=A_R\mathbf e_R+A_\phi\mathbf e_\phi+A_Z\mathbf e_Z$,
\begin{equation}
\nabla\cdot\mathbf A
=\frac{1}{R}\pdv{}{R}(RA_R)
+\frac{1}{R}\pdv{A_\phi}{\phi}
+\pdv{A_Z}{Z}.
\end{equation}
The $R$ factor encodes geometry, not a correction that may be dropped.

\begin{physicsbox}
A conservative discretization should imitate the integral theorem behind the differential operator. Discrete divergence should be net face flux divided by cell volume.
\end{physicsbox}


\section{Computational laboratory}

On a Cartesian grid, define an analytic vector field and compute its divergence and curl exactly and numerically. Demonstrate second-order convergence. Verify the divergence theorem by comparing volume integral and boundary flux. Repeat for an axisymmetric cylindrical field and show the error from omitting the $1/R$ factor.


\section{Summary}
\begin{itemize}[leftmargin=1.6em]
\item Gradient measures directional scalar change.
\item Divergence is the local operator of conservation laws.
\item Curl measures circulation and connects currents and induction.
\item Vector identities reveal magnetic pressure, tension, stretching, and advection.
\item Geometry must be built into both analytic and discrete operators.
\end{itemize}

\SymbolsAcronymsSection
\begin{symtable}
$\nabla f$ & gradient & Steepest scalar increase. \\
$\nabla\cdot\mathbf F$ & divergence & Net outward flux per volume. \\
$\nabla\times\mathbf F$ & curl & Local circulation density. \\
$\nabla^2$ & Laplacian & Divergence of gradient. \\
$d\mathbf S$ & oriented area & Normal direction times surface area. \\
\end{symtable}

\section{Exercises}
\begin{enumerate}[leftmargin=1.8em]
\item Use the divergence theorem with $\mathbf F=\mathbf r/3$ to derive the volume of a sphere.
\item Prove $\nabla\cdot(\nabla\times\mathbf A)=0$ for smooth Cartesian fields.
\item Derive the magnetic pressure-tension identity using index notation.
\item Compute divergence and curl of rigid rotation $\mathbf u=(-y,x,0)$.
\item Write the integral form of the continuity equation for a control volume.
\end{enumerate}

\section{References}
\begin{thebibliography}{99}
\bibitem{ref1} G. B. Arfken, H. J. Weber, and F. E. Harris, \emph{Mathematical Methods for Physicists}, 7th ed., Academic Press (2013).
\bibitem{ref2} G. Strang, \emph{Linear Algebra and Its Applications}, 4th ed., Brooks/Cole (2006).
\bibitem{ref3} W. A. Strauss, \emph{Partial Differential Equations: An Introduction}, 2nd ed., Wiley (2007).
\bibitem{ref4} D. J. Griffiths, \emph{Introduction to Electrodynamics}, 4th ed., Cambridge University Press (2017).
\bibitem{ref5} J. D. Jackson, \emph{Classical Electrodynamics}, 3rd ed., Wiley (1998).
\bibitem{ref6} J. P. Freidberg, \emph{Ideal MHD}, Cambridge University Press (2014).
\bibitem{ref7} J. P. Goedbloed, R. Keppens, and S. Poedts, \emph{Advanced Magnetohydrodynamics}, Cambridge University Press (2010).
\bibitem{ref8} H. Goedbloed and S. Poedts, \emph{Principles of Magnetohydrodynamics}, Cambridge University Press (2004).
\end{thebibliography}

\end{document}
